\documentclass[11pt]{article}

\usepackage[margin=0.85in]{geometry}
\usepackage[T1]{fontenc}
\usepackage[utf8]{inputenc}
\usepackage{lmodern}
\usepackage{microtype}
\usepackage{hyperref}

\hypersetup{
  colorlinks=true,
  linkcolor=blue,
  urlcolor=blue,
  citecolor=blue
}

\setlength{\parskip}{0.55em}
\setlength{\parindent}{0pt}

\title{Cache-Oriented Performance Opportunities in Pioneer.jl}
\author{Prepared for Pioneer.jl development}
\date{May 18, 2026}

\begin{document}
\maketitle

\section*{Context}

This note maps the main lessons from Algorithmica's ``RAM \& CPU Caches''
chapter to the current Pioneer.jl codebase. The relevant cache concepts are:
bandwidth vs. latency, cache-line traffic, memory-level parallelism, hardware
prefetching, structure layout, cache associativity, TLB/page behavior, and
AoS-vs-SoA layout.

Pioneer already uses several of these ideas in its performance-critical code:
compact fragment records, struct-of-arrays fragment-bin metadata, partition-local
\path{UInt16} precursor IDs, partition-major traversal, fused match/build loops,
and packed metadata in the fused sparse matrix. The best opportunities now are
not broad rewrites. They are targeted reductions in avoidable memory traffic,
temporary allocation, and pointer chasing.

\section*{Priority Summary}

\begin{itemize}
  \item \textbf{Precompute top-N intensity thresholds once per scan.}
        Low risk. Removes repeated scans and partial sorts when one scan
        overlaps multiple fragment-index partitions.
  \item \textbf{Avoid per-scan candidate allocation in fused search.}
        Low risk. Removes hot-loop allocation and copy work while preserving the
        existing algorithm.
  \item \textbf{Repack PSM accumulator structs.}
        Medium risk. Hot records are currently about 72 bytes, just over one
        64-byte cache line. Better packing can reduce cache-line pressure.
  \item \textbf{Benchmark dense or generation-stamped precursor maps.}
        Medium risk. Current \path{Dict}-backed maps avoid large per-thread
        arrays but introduce pointer chasing in scan scoring.
  \item \textbf{Specialize simple mass-error peak matching.}
        Medium risk. The general path loads three peak-window arrays; simple
        models may need fewer streams.
  \item \textbf{Use fixed-size isotope scratch storage.}
        Medium risk. Tiny vector scratch is repeatedly filled and read inside
        the fragment inner loop.
\end{itemize}

\section{Precompute Top-N Intensity Thresholds Once Per Scan}

\subsection*{Where}

The partitioned fragment-index search computes an optional top-N peak intensity
threshold inside \path{_run_thread}:

\begin{itemize}
  \item \path{src/structs/SpectralLibrary/PartitionedFragmentIndex/search.jl:519}
  \item threshold extraction around \path{search.jl:552}
  \item \path{partialsort!} call around \path{search.jl:559}
\end{itemize}

This sits inside the partition-major loop. A single scan can be relevant to
more than one precursor-m/z partition. When that happens, its intensity array is
visited and partially sorted once per relevant partition.

\subsection*{Cache concept}

This is a cache-line bandwidth issue. The intensity array is contiguous and
therefore friendly to hardware prefetching, but repeatedly rereading it still
consumes memory bandwidth and displaces other hot data. The \path{partialsort!}
also mutates a per-thread scratch buffer and can touch a substantial fraction of
the scan's peak intensities.

\subsection*{Implementation sketch}

Move the threshold calculation into the existing per-scan precompute phase in
\path{searchFragmentIndexPartitionMajorHinted}. The function already computes:

\begin{itemize}
  \item scan iRT lower/upper bounds,
  \item scan precursor m/z bounds,
  \item scan center iRT.
\end{itemize}

Add a \path{Vector{Float32}} named something like \path{scan_intensity_thresholds}.
If \path{max_peaks == 0}, fill it with \path{0.0f0}. Otherwise, compute the
threshold once per \path{all_scan_idxs} entry before the partition-major
parallel loop. Then pass this vector into \path{_run_thread} and read:

\begin{verbatim}
intensity_threshold = scan_intensity_thresholds[si]
\end{verbatim}

inside the scan loop.

\subsection*{Expected benefit}

The win should scale with:

\begin{itemize}
  \item the number of scans that overlap multiple partitions,
  \item the number of peaks per scan,
  \item the number of active threads,
  \item whether \path{max_peaks > 0}.
\end{itemize}

This should reduce both runtime and transient memory traffic during wide scout
or other top-N filtered index searches. The behavior should be exactly
equivalent if the threshold definition is preserved.

\subsection*{Risks}

The main risk is memory overhead for storing one threshold per scan, which is
small. A secondary risk is moving work from per-thread scratch into a shared
precompute stage. That should be fine because each scan threshold is independent.

\subsection*{Benchmark}

Use a representative file and compare:

\begin{itemize}
  \item total time in \path{searchFragmentIndexPartitionMajorHinted},
  \item allocations,
  \item number of threshold computations, if instrumented,
  \item candidate count equality.
\end{itemize}

The candidate output should be bit-identical because only redundant threshold
calculation is removed.

\section{Avoid Per-Scan Candidate Allocation in Fused Search}

\subsection*{Where}

In the fused scan path:

\begin{itemize}
  \item \path{src/Routines/SearchDIA/process_scans_fused.jl:21}
  \item \path{collect(Int64, prec_range)} around \path{process_scans_fused.jl:146}
  \item optional target/decoy split buffers around \path{process_scans_fused.jl:134}
\end{itemize}

The normal non-split path materializes a fresh \path{Vector{Int64}} from
\path{prec_range} before calling \path{run_fused!}. But \path{run_fused!} already
accepts an \path{AbstractVector{Int64}} for \path{prec_range}, and the loop body
does not require ownership of the candidate index list.

\subsection*{Cache concept}

This combines bandwidth and allocation overhead. Allocating and filling a vector
per scan turns an already available range into fresh memory traffic. It also adds
pressure on Julia's allocator and garbage collector.

\subsection*{Implementation sketch}

There are two practical routes:

\begin{enumerate}
  \item Broaden \path{run_fused!}'s \path{prec_range} type to accept
        \path{AbstractUnitRange{<:Integer}} or simply remove the concrete
        \path{AbstractVector{Int64}} restriction.
  \item Call \path{_run_subset!(prec_range)} directly in the non-split branch.
\end{enumerate}

For the split branch, keep two reusable buffers on the per-thread search context
or fused scratch object, for example \path{target_idx_scratch} and
\path{decoy_idx_scratch}. Use \path{empty!} between scans. That keeps the split
feature allocation-free after warmup.

\subsection*{Expected benefit}

This is likely a small but reliable win. It should show up most clearly in
files with many MS2 scans and modest candidate counts where allocation overhead
is not drowned out by deconvolution time.

\subsection*{Risks}

Low. The non-split path should preserve behavior exactly. For split mode, the
main risk is accidentally retaining stale entries if buffers are not cleared.
That is easy to test.

\subsection*{Benchmark}

Run the current unit tests around fused scanning, then compare allocation counts
for a representative search with:

\begin{verbatim}
@time
@allocated
\end{verbatim}

or the existing profiling hooks. Confirm identical PSM output for a fixed input.

\section{Repack PSM Accumulator and Output Structs}

\subsection*{Where}

The hot PSM records are:

\begin{itemize}
  \item \path{MainUnscoredPSM{Float32}} in \path{src/Routines/SearchDIA/PSMs/UnscoredPSMs.jl:27}
  \item \path{TuningUnscoredPSM{Float32}} in the same file
  \item \path{MainSearchScoredPSM{Float32,Float16}} in \path{src/Routines/SearchDIA/PSMs/ScoredPSMs.jl:70}
  \item \path{TuningScoredPSM{Float32,Float16}} in the same file
\end{itemize}

A quick size check showed:

\begin{itemize}
  \item \path{MainUnscoredPSM{Float32}}: 72 bytes,
  \item \path{TuningUnscoredPSM{Float32}}: 36 bytes,
  \item \path{MainSearchScoredPSM{Float32,Float16}}: 72 bytes,
  \item \path{TuningScoredPSM{Float32,Float16}}: 40 bytes.
\end{itemize}

The two main-search structs cross the 64-byte cache-line boundary.

\subsection*{Cache concept}

CPU caches move data in cache lines, commonly 64 bytes. A 72-byte record often
requires two cache lines when accessed in an array, even if the code only needs
part of the record. That increases bandwidth use and lowers effective cache
capacity.

\subsection*{Implementation sketch}

First, reorder fields to group by alignment:

\begin{itemize}
  \item all \path{Float32} fields together,
  \item all \path{Float16} fields together,
  \item all \path{UInt32} fields together,
  \item all \path{UInt8} fields together.
\end{itemize}

Then consider packing low-cardinality counters:

\begin{itemize}
  \item several \path{UInt8} counts can be packed into one \path{UInt32},
  \item boolean/rank flags can be packed into bit fields represented by
        \path{UInt16} or \path{UInt32},
  \item top-rank intensities should stay as separate \path{Float32} fields if
        downstream code reads them directly.
\end{itemize}

This should be done with constructors and accessor functions so the scoring code
does not become unreadable. The current \path{CompactFrag} design is a good local
pattern: it packs metadata but keeps inline getters.

\subsection*{Expected benefit}

If \path{MainUnscoredPSM} can be brought to 64 bytes or less, updates in
\path{apply_main_scoring!} and reads in \path{Score!} should touch fewer cache
lines. If \path{MainSearchScoredPSM} can also be reduced, output construction
and later ML feature-table conversion may benefit.

\subsection*{Risks}

This is riskier than the first two ideas because these structs are part of
internal data contracts. Field order changes can affect Arrow output if any code
relies on Tables.jl's struct field order. Packing also adds getter/setter logic,
which can make code harder to maintain.

The safe path is:

\begin{enumerate}
  \item add size tests using \path{sizeof},
  \item preserve public column names in output conversion,
  \item benchmark before and after,
  \item avoid packing fields that are read constantly unless the getter inlines
        cleanly.
\end{enumerate}

\subsection*{Benchmark}

Benchmark:

\begin{itemize}
  \item fused scan runtime,
  \item \path{Score!} runtime,
  \item PSM output equality,
  \item \path{sizeof} for each struct,
  \item allocations and compilation impact.
\end{itemize}

\section{Benchmark Dense or Generation-Stamped Precursor Maps}

\subsection*{Where}

The scan-local precursor map abstraction is in:

\begin{itemize}
  \item \path{src/structs/PrecursorMap.jl}
  \item dense backend around \path{PrecursorMap.jl:57}
  \item sparse \path{Dict}-backed backend around \path{PrecursorMap.jl:159}
  \item current search context field \path{id_to_col::SparsePrecMap{UInt16}}
        in \path{src/Routines/SearchDIA/SearchMethods/SearchTypes.jl:196}
\end{itemize}

\subsection*{Cache concept}

Julia \path{Dict} lookup is flexible and memory-efficient for sparse active
sets, but it involves hashing and indirect memory access. That can introduce
pointer chasing, which is exactly the kind of access pattern that defeats memory
prefetching and exposes memory latency.

A dense vector indexed by precursor ID has excellent locality for each lookup,
but it can be large per thread. The repo comments already note that a fully
dense map can cost substantial memory on large libraries.

\subsection*{Implementation sketch}

Benchmark three variants behind the existing \path{AbstractPrecursorMap}
interface:

\begin{enumerate}
  \item current \path{SparsePrecMap};
  \item existing \path{DensePrecMap};
  \item generation-stamped dense map.
\end{enumerate}

A generation-stamped dense map would store:

\begin{itemize}
  \item \path{vals::Vector{V}},
  \item \path{stamp::Vector{UInt16}} or \path{Vector{UInt32}},
  \item current generation counter,
  \item active key list.
\end{itemize}

Reset becomes:

\begin{verbatim}
generation += 1
n_active = 0
\end{verbatim}

Instead of zeroing active entries. A key is present if its stamp equals the
current generation. This is useful when the dense vector is acceptable but reset
or active tracking should be cheap.

\subsection*{Expected benefit}

This could speed up:

\begin{itemize}
  \item \path{id_to_col[precursor_idx]} in \path{Score!},
  \item \path{id_to_col[match_column_id(...)] = this_col} in \path{run_fused!},
  \item warm-start weight transfer in deconvolution utilities.
\end{itemize}

The expected benefit is workload-dependent. For very large libraries and many
threads, the dense memory footprint may hurt more than it helps. For moderate
libraries or machines with enough memory, it may remove a latency-heavy hash
lookup from hot paths.

\subsection*{Risks}

The main risk is per-thread memory growth and lower effective cache capacity.
This is exactly the tradeoff highlighted by the cache chapter: a theoretically
faster direct lookup can become slower if it forces a much larger working set.

\subsection*{Benchmark}

Make backend selection configurable by an environment variable or parameter,
then compare:

\begin{itemize}
  \item fused search runtime,
  \item max resident memory,
  \item GC time,
  \item output equality.
\end{itemize}

The result should be interpreted by library size and thread count, not as a
single universal winner.

\section{Specialize Simple Mass-Error Peak Matching}

\subsection*{Where}

The fused path precomputes peak matching windows in:

\begin{itemize}
  \item \path{prepare_scan_peaks!} at \path{src/Routines/SearchDIA/CommonSearchUtils/fusedMatch.jl:417}
  \item SIMD matching in \path{scan_for_nearest_in_window} at \path{fusedMatch.jl:476}
\end{itemize}

The current general path stores and scans three contiguous arrays:

\begin{itemize}
  \item corrected peak m/z,
  \item lower observable bound,
  \item upper observable bound.
\end{itemize}

This is a good SoA design for the general intensity-aware mass-error model.

\subsection*{Cache concept}

SoA helps SIMD because each vector load reads one field across adjacent peaks.
But every extra array is another memory stream. The current matching loop can
load three 32-byte vectors for every eight peaks. If a simpler mass-error model
does not need precomputed per-peak bounds, one or two of those streams may be
avoidable.

\subsection*{Implementation sketch}

Keep the general path unchanged. Add a specialized method for simple mass-error
models where the tolerance can be computed directly from the theoretical m/z or
corrected m/z without reading \path{obs_low} and \path{obs_high}.

For example:

\begin{itemize}
  \item for a fixed ppm model, search by corrected m/z and compare absolute
        error to \path{iso_mz * ppm * 1e-6};
  \item for a constant Da model, compare absolute error to the constant bound;
  \item keep the intensity-aware model on the existing three-array path.
\end{itemize}

This would be multiple dispatch on the concrete mass-error model type, not a
runtime branch inside the hottest loop.

\subsection*{Expected benefit}

For simple models, this can reduce peak-window scanning bandwidth by up to two
array streams. It may also reduce the per-scan precompute work in
\path{prepare_scan_peaks!}. The effect should be largest on scans with many
peaks and many candidate fragment isotope windows.

\subsection*{Risks}

The semantic risk is high if the specialized path does not exactly match the
mass-error model's intended acceptance window. The current code explicitly
preserves intensity-dependent bias and tolerance. That behavior must stay
untouched for those models.

The implementation should therefore start with a very narrow specialization for
one simple model and an equivalence test against the current path.

\subsection*{Benchmark}

For a fixed simple model:

\begin{itemize}
  \item compare matched peak indices and ppm errors against the current path,
  \item measure fused matching time,
  \item inspect allocation and cache behavior if hardware counters are available.
\end{itemize}

\section{Use Fixed-Size Isotope Scratch Storage}

\subsection*{Where}

Fragment isotope intensities are computed in:

\begin{itemize}
  \item \path{src/Routines/SearchDIA/CommonSearchUtils/getFragIsotopes.jl:17}
  \item called inside \path{run_fused!} around \path{src/Routines/SearchDIA/CommonSearchUtils/fusedMatch.jl:818}
\end{itemize}

The search context currently stores \path{isotopes::Vector{Float32}} and
\path{precursor_transmission::Vector{Float32}}.

\subsection*{Cache concept}

The isotope buffers are tiny, but they sit in the innermost fragment loop. The
current code repeatedly fills and reads a heap-allocated vector. That is not a
large memory object, but the operation is frequent. Fixed-size storage can keep
more of the state in registers or stack-like local storage, improving locality
and reducing bounds checks when the compiler can see the size.

\subsection*{Implementation sketch}

There are two levels of aggressiveness:

\begin{enumerate}
  \item Replace \path{fill!(frag_isotopes, 0)} with loops only over the active
        isotope range where safe. This is the conservative option.
  \item Introduce a fixed-size representation for the common maximum isotope
        count, for example an \path{NTuple} or \path{MVector} style buffer if
        StaticArrays is already acceptable in the hot path.
\end{enumerate}

The second route needs care. StaticArrays can help for tiny fixed-size data, but
using them incorrectly can increase register pressure or compilation time.

\subsection*{Expected benefit}

This is likely a small inner-loop win rather than a large algorithmic gain. It
may matter because \path{getFragIsotopes!} is called for every candidate
fragment that passes rank filters.

\subsection*{Risks}

The main risk is increasing code complexity in a numerically sensitive area.
Another risk is creating too many method specializations if isotope count
becomes a type parameter across search kinds.

\subsection*{Benchmark}

Microbenchmark \path{getFragIsotopes!} for common fragment and precursor cases,
then benchmark full fused search. Require bitwise or near-bitwise equality for
isotope intensities and PSM results.

\section*{General Benchmarking Plan}

For each candidate, use a two-layer benchmark:

\begin{enumerate}
  \item A focused microbenchmark for the changed function.
  \item An end-to-end representative search file.
\end{enumerate}

Track:

\begin{itemize}
  \item wall time,
  \item allocations,
  \item output equality or accepted numerical tolerance,
  \item resident memory,
  \item number of candidates, matches, misses, and PSMs.
\end{itemize}

If available, hardware counters would be ideal for the cache-specific claims:
L1/L2/L3 misses, LLC load misses, branch misses, and cycles. Without counters,
allocation counts and stable end-to-end timing are still enough to decide
whether a change is worth keeping.

\section*{Suggested Order of Work}

\begin{enumerate}
  \item Implement top-N threshold precompute. It is low risk and likely
        measurable.
  \item Remove fused per-scan candidate allocation. This is straightforward and
        should reduce allocation noise.
  \item Add a backend switch for precursor maps and benchmark sparse vs. dense
        vs. generation-stamped dense.
  \item Repack PSM structs only after the first two changes are measured.
  \item Try simple-model peak matching specialization if profiling shows peak
        matching remains dominant.
  \item Try fixed-size isotope scratch last, because it is the most inner-loop
        specific and easiest to overfit.
\end{enumerate}

\section*{References}

\begin{itemize}
  \item Algorithmica, RAM \& CPU Caches:
        \url{https://en.algorithmica.org/hpc/cpu-cache/}
  \item Memory bandwidth:
        \url{https://en.algorithmica.org/hpc/cpu-cache/bandwidth/}
  \item Memory latency:
        \url{https://en.algorithmica.org/hpc/cpu-cache/latency/}
  \item Cache lines:
        \url{https://en.algorithmica.org/hpc/cpu-cache/cache-lines/}
  \item Memory-level parallelism and prefetching:
        \url{https://en.algorithmica.org/hpc/cpu-cache/mlp/},
        \url{https://en.algorithmica.org/hpc/cpu-cache/prefetching/}
  \item Alignment, associativity, paging, and AoS/SoA:
        \url{https://en.algorithmica.org/hpc/cpu-cache/alignment/},
        \url{https://en.algorithmica.org/hpc/cpu-cache/associativity/},
        \url{https://en.algorithmica.org/hpc/cpu-cache/paging/},
        \url{https://en.algorithmica.org/hpc/cpu-cache/aos-soa/}
\end{itemize}

\end{document}
